
\documentclass[11pt,a4paper]{article}

\usepackage{fontspec}
\usepackage[french]{babel}
\usepackage{amsmath,amssymb,amsthm,mathtools}
\usepackage{geometry}
\usepackage{xcolor}
\usepackage{fancyhdr}
\usepackage{enumitem}
\usepackage{hyperref}
\usepackage{titlesec}
\usepackage{microtype}
\usepackage{booktabs}
\usepackage{array}

\geometry{margin=2.2cm}
\setmainfont{DejaVu Serif}
\setsansfont{DejaVu Sans}
\setmonofont{DejaVu Sans Mono}

\definecolor{accent}{HTML}{1F4E79}
\definecolor{lightaccent}{HTML}{EAF2F8}

\pagestyle{fancy}
\fancyhf{}
\lhead{\small Cours Deep Reinforcement Learning}
\chead{\small Classe : 4IA}
\rhead{\small Année : 2026}
\lfoot{\small Enseignant : Mourad Zéraï, PhD - ESPRIT School of Engineering}
\cfoot{}
\rfoot{\thepage}
\setlength{\headheight}{14pt}

\titleformat{\section}{\large\bfseries\color{accent}}{\thesection.}{0.5em}{}
\titleformat{\subsection}{\normalsize\bfseries\color{accent}}{\thesubsection.}{0.5em}{}
\titleformat{\subsubsection}{\normalsize\bfseries}{\thesubsubsection.}{0.5em}{}

\newtheorem*{definition}{Définition}
\newtheorem*{rappel}{Rappel}
\newtheorem*{hypothese}{Hypothèse}
\newtheorem*{propriete}{Propriété}
\newtheorem*{theoreme}{Théorème}
\newtheorem*{remarque}{Remarque}

\newcommand{\E}{\mathbb{E}}
\newcommand{\Pbb}{\mathbb{P}}
\newcommand{\R}{\mathbb{R}}
\newcommand{\argmax}{\operatorname*{argmax}}
\newcommand{\given}{\,|\,}

\begin{document}

\begin{center}
    {\LARGE\bfseries Dérivation et explication pas à pas de Policy Iteration}\\[0.4em]
    {\large Note de cours rigoureuse, détaillée et appliquée à FrozenLake}
\end{center}

\vspace{0.6em}

\section{Objectif}

Le but de cette note est d'expliquer rigoureusement \emph{Policy Iteration}, d'en justifier les deux briques mathématiques - évaluation de politique et amélioration de politique - puis d'en montrer l'application à FrozenLake.

\section{Cadre et définitions}

On considère un MDP fini avec $\gamma\in[0,1)$.

\begin{definition}[Valeur d'une politique]
Pour une politique $\pi$,
\[
v_{\pi}(s)=\E_{\pi}[G_t \given S_t=s].
\]
\end{definition}

\begin{definition}[Valeur état-action associée à $\pi$]
\[
q_{\pi}(s,a)=\sum_{s',r} p(s',r\given s,a)\left[r+\gamma v_{\pi}(s')\right].
\]
\end{definition}

\begin{definition}[Politique gloutonne par rapport à $v$]
Une politique $\pi'$ est dite gloutonne par rapport à une fonction $v$ si, pour tout état $s$,
\[
\pi'(s)\in \argmax_a \sum_{s',r} p(s',r\given s,a)\left[r+\gamma v(s')\right].
\]
\end{definition}

Policy Iteration alterne :
\begin{enumerate}[leftmargin=1.6em]
    \item une phase d'évaluation exacte de la politique courante ;
    \item une phase d'amélioration gloutonne.
\end{enumerate}

\section{Étape 1 : évaluation de politique}

\subsection{Équation vérifiée par $v_{\pi}$}

Pour une politique fixée $\pi$,
\[
v_{\pi}(s)=\sum_a \pi(a\given s)\sum_{s',r} p(s',r\given s,a)\left[r+\gamma v_{\pi}(s')\right].
\]
C'est l'équation de Bellman d'espérance.

\subsection{Forme opérateur}

On définit l'opérateur
\[
(T_{\pi}v)(s)=\sum_a \pi(a\given s)\sum_{s',r} p(s',r\given s,a)\left[r+\gamma v(s')\right].
\]
Alors
\[
v_{\pi}=T_{\pi}v_{\pi}.
\]

\subsection{Justification de la convergence de l'évaluation}

Comme pour Value Iteration, on montre que $T_{\pi}$ est une contraction.

\begin{theoreme}
Pour tous $v,w$,
\[
\|T_{\pi}v-T_{\pi}w\|_{\infty}\le \gamma\|v-w\|_{\infty}.
\]
\end{theoreme}

\paragraph{Justification.}
Pour un état $s$,
\begin{align*}
|(T_{\pi}v)(s)-(T_{\pi}w)(s)|
&= \left|\sum_a \pi(a\given s)\sum_{s',r}p(s',r\given s,a)\gamma\left[v(s')-w(s')\right]\right| \\
&\le \sum_a \pi(a\given s)\sum_{s',r}p(s',r\given s,a)\gamma |v(s')-w(s')| \\
&\le \gamma \|v-w\|_{\infty}\sum_a \pi(a\given s)\sum_{s',r}p(s',r\given s,a).
\end{align*}
Comme $\sum_a \pi(a\given s)=1$ et $\sum_{s',r}p(s',r\given s,a)=1$, on obtient
\[
|(T_{\pi}v)(s)-(T_{\pi}w)(s)|\le \gamma\|v-w\|_{\infty}.
\]
Puis on prend le maximum en $s$.

Donc l'itération
\[
v^{(k+1)}=T_{\pi}v^{(k)}
\]
converge vers l'unique point fixe $v_{\pi}$.

\section{Étape 2 : amélioration de politique}

\subsection{Principe}

Une fois $v_{\pi}$ calculée, on définit une nouvelle politique $\pi'$ gloutonne :
\[
\pi'(s)\in\argmax_a \sum_{s',r}p(s',r\given s,a)\left[r+\gamma v_{\pi}(s')\right].
\]

\subsection{Pourquoi cette politique est meilleure}

Le résultat central est le théorème d'amélioration de politique.

\begin{theoreme}[Policy Improvement]
Si, pour tout état $s$,
\[
q_{\pi}(s,\pi'(s)) \ge v_{\pi}(s),
\]
alors
\[
v_{\pi'}(s)\ge v_{\pi}(s)\qquad \text{pour tout } s.
\]
\end{theoreme}

\paragraph{Justification pas à pas.}
Comme $\pi'$ est gloutonne par rapport à $v_{\pi}$,
\[
q_{\pi}(s,\pi'(s)) = \max_a q_{\pi}(s,a)\ge q_{\pi}(s,\pi(s))=v_{\pi}(s).
\]
Le dernier passage vient de l'identité
\[
v_{\pi}(s)=q_{\pi}(s,\pi(s))
\]
dans le cas d'une politique déterministe. Plus généralement,
\[
v_{\pi}(s)=\sum_a \pi(a\given s)q_{\pi}(s,a).
\]

On a donc
\[
v_{\pi}(s)\le \E_{\pi'}[R_{t+1}+\gamma v_{\pi}(S_{t+1})\given S_t=s].
\]
On réinjecte cette inégalité à l'instant suivant, puis encore, puis encore. On obtient une suite d'inégalités emboîtées qui, après itération, donne
\[
v_{\pi}(s)\le \E_{\pi'}\!\left[\sum_{k=0}^{n-1}\gamma^k R_{t+k+1}+\gamma^n v_{\pi}(S_{t+n}) \middle\vert S_t=s\right].
\]
Sous l'hypothèse $\gamma<1$ et avec des récompenses bornées, le terme résiduel $\gamma^n v_{\pi}(S_{t+n})$ tend vers $0$ lorsque $n\to\infty$. En passant à la limite, on obtient
\[
v_{\pi}(s)\le \E_{\pi'}[G_t\given S_t=s]=v_{\pi'}(s).
\]

\section{Algorithme complet}

\subsection{Initialisation}

On choisit une politique initiale $\pi_0$.

\subsection{Boucle principale}

À l'itération $k$ :

\begin{enumerate}[leftmargin=1.6em]
    \item \textbf{Évaluation} : calculer exactement ou numériquement $v_{\pi_k}$.
    \item \textbf{Amélioration} : définir
    \[
    \pi_{k+1}(s)\in \argmax_a \sum_{s',r} p(s',r\given s,a)\left[r+\gamma v_{\pi_k}(s')\right].
    \]
\end{enumerate}

\subsection{Critère d'arrêt}

Si $\pi_{k+1}=\pi_k$, alors la politique est stable. Dans un MDP fini, une politique stable est optimale.

\section{Pourquoi une politique stable est optimale}

Si $\pi$ est stable, alors pour tout état
\[
\pi(s)\in \argmax_a \sum_{s',r}p(s',r\given s,a)\left[r+\gamma v_{\pi}(s')\right].
\]
Donc
\[
v_{\pi}(s)=\max_a \sum_{s',r}p(s',r\given s,a)\left[r+\gamma v_{\pi}(s')\right].
\]
Autrement dit, $v_{\pi}$ satisfait l'équation d'optimalité de Bellman. Par unicité du point fixe de l'opérateur optimal, on en déduit
\[
v_{\pi}=v_*.
\]
Donc $\pi$ est optimale.

\section{Application à FrozenLake}

\subsection{Politique initiale}

Dans FrozenLake 4x4, on peut initialiser par exemple une politique uniforme :
\[
\pi_0(a\given s)=\frac14
\]
sur les états non terminaux, ou une politique déterministe arbitraire.

\subsection{Évaluation dans FrozenLake}

Comme le modèle de transition est connu via \texttt{env.unwrapped.P}, on peut évaluer exactement la politique par itération de Bellman :
\[
v^{(k+1)}(s)=\sum_a \pi(a\given s)\sum_i p_i\left[r_i+\gamma v^{(k)}(s'_i)\right].
\]

\subsection{Amélioration dans FrozenLake}

Pour chaque case $s$ et action $a$ :
\[
q_{\pi}(s,a)=\sum_i p_i\left[r_i+\gamma v_{\pi}(s'_i)\right].
\]
Puis on choisit l'action de plus grande valeur.

\subsection{Effet du mode glissant}

En mode glissant, une action n'entraîne pas toujours la case désirée. L'amélioration de politique tient compte de ce risque. Une case proche d'un trou peut devenir moins attractive si une glissade y mène souvent.

\section{Pseudo-code pédagogique}

\begin{verbatim}
Choisir une politique initiale pi

Répéter :
    Evaluer pi jusqu'à convergence pour obtenir v_pi

    policy_stable = True
    pour chaque état s :
        old_action = pi(s)
        pi(s) = argmax_a somme_{s',r} p(s',r|s,a)[r + gamma v_pi(s')]
        si pi(s) != old_action :
            policy_stable = False
Jusqu'à ce que policy_stable soit vraie
\end{verbatim}

\section{Ce qu'il faut retenir}

\begin{itemize}[leftmargin=1.6em]
    \item Policy Iteration sépare clairement \emph{évaluation} et \emph{amélioration}.
    \item L'évaluation repose sur l'opérateur $T_{\pi}$, qui est une contraction.
    \item L'amélioration gloutonne ne détériore jamais la politique et l'améliore souvent strictement.
    \item Une politique stable est optimale.
    \item Dans FrozenLake, l'algorithme est particulièrement transparent car le MDP est fini et le modèle est explicite.
\end{itemize}

\end{document}
